\section{Situación Actual}

\subsection{Descripción del problema}

En la era de la hiperconectividad, el entorno digital se ha transformado en un arma de doble filo para la infancia. La rápida expansión del acceso a internet y de las tecnologías digitales ha transformado profundamente la vida de millones de niños y adolescentes alrededor del mundo. Actualmente, más de dos tercios de los menores de edad utilizan internet de forma habitual, ya sea para fines educativos, recreativos o sociales.\footnote{UNICEF. (s. f.). El Estado Mundial de la Infancia 2017: Los niños en un mundo digital — Resumen. \url{https://www.unicef.org/media/48611/file/SOWC_2017_Summary_SP.pdf}} Si bien es cierto que esto ha generado importantes oportunidades para el aprendizaje y la inclusión digital, el acceso universal a internet y la proliferación de plataformas digitales también han expuesto a los menores a riesgos sin precedentes como el grooming, sextorsión, la producción y distribución de material de abuso sexual infantil, la trata de personas y el reclutamiento por parte de grupos extremistas o criminales.

Las plataformas digitales han reducido considerablemente las barreras para que estos actores contacten a menores de cualquier parte del mundo. Redes sociales, videojuegos en línea, plataformas de transmisión en vivo e incluso herramientas educativas han sido utilizadas para establecer contacto con posibles víctimas. Organizaciones como el National Center for Missing and Exploited Children han reportado un incremento exponencial en las denuncias de material de abuso sexual infantil en los últimos años; en 2025, la cifra de reportes emitidos fue de 21.3 millones.\footnote{Missing \& Exploited Children. (s. f.). Datos de la CyberTipline. Centro Nacional para Niños Desaparecidos y Explotados. \url{https://www.missingkids.org/gethelpnow/cybertipline/cybertiplinedata}} Adicionalmente, el fenómeno de la sextorsión financierase ha disparado globalmente, afectando de manera desproporcionada tanto a adolescentes vulnerables como a niños de rangos de edad cada vez menores.\footnote{Lisa Institute. (s. f.). Grooming, sexting, stalking y sextorsión: Definición y diferencias. \url{https://www.lisainstitute.com/blogs/blog/grooming-sexting-stalking-sextorsion-definicion-diferencias}} \footnote{Universidad de Edimburgo. (s. f.). Un informe muestra la magnitud del abuso infantil en Internet. \url{https://www.ed.ac.uk/news/report-shows-scale-of-online-abuse-of-children}} Asociaciones internacionales han advertido que este tipo de delito ha aumentado considerablemente desde la pandemia de COVID-19 debido al incremento del tiempo que los menores pasan conectados.

Simultáneamente, diversos grupos extremistas continúan utilizando internet como herramienta para la radicalización de menores.\footnote{The Soufan Center. (s. f.). Lucha contra la radicalización de los jóvenes. \url{https://thesoufancenter.org/research/combating-youth-radicalization/}} \footnote{UNICEF. (s. f.). Cómo mantener a los niños seguros en Internet. \url{https://www.unicef.org/protection/keeping-children-safe-online}} Aprovechándose de la vulnerabilidad emocional, estos grupos logran difundir propaganda, crear comunidades virtuales y atraer nuevos integrantes a través de diferentes páginas de internet. La dimensión transnacional del problema representa uno de los principales desafíos para los estados. Un delito puede involucrar a un agresor ubicado en un país, una víctima en otro, servidores alojados en un tercero y plataformas digitales cuya sede legal se encuentra en una jurisdicción completamente distinta.

A su vez, el desarrollo acelerado de la inteligencia artificial ha generado una nueva dimensión al problema.\footnote{World Economic Forum. (2025, enero). Avances tecnológicos y desarrollo humano: Una historia de dos mundos. \url{https://es.weforum.org/stories/2025/01/avances-tecnologicos-y-desarrollo-humano-una-historia-de-dos-m} undos/} Herramientas capaces de generar imágenes hiperrealistas, modificar voces o producir videos falsificados (deep fakes) han comenzado a ser utilizadas para crear material sexual que involucra menores o para facilitar la manipulación y el engaño. Aunque algunos de estos contenidos no implican a víctimas reales, su existencia plantea importantes desafíos legales, éticos y de protección infantil. Sin embargo, la inteligencia artificial también ofrece nuevas oportunidades para combatir estos delitos.

Sistemas automatizados de reconocimiento de imágenes, análisis de patrones de comportamiento y detección de redes criminales permiten identificar con mayor rapidez contenido ilegal y apoyar las investigaciones.

\subsection{Planteamiento del problema}

La persistencia y el agravamiento de esta problemática no se deben solo a una causa, sino a una compleja red de factores de vulnerabilidad individuales y fallas estructurales a nivel sistemático y regulatorio. Suele existir una combinación de condiciones sociales, psicológicas, tecnológicas e institucionales que juntas explican esta victimización.

Uno de los principales factores de riesgo es la limitada alfabetización digital.\footnote{UNESCO. (s. f.). Educación digital. \url{https://www.unesco.org/en/digital-education}} Aunque las generaciones más jóvenes utilizan internet desde edades tempranas, ello no implica necesariamente que comprendan los riesgos asociados a la interacción con desconocidos, la protección de datos personales o la identificación de mecanismos de manipulación utilizados por delincuentes y grupos extremistas.\footnote{Organización de las Naciones Unidas para la Educación, la Ciencia y la Cultura. (s. f.). Preguntas y respuestas: ¿Por qué es esencial la educación digital para la ciudadanía mundial? \url{https://www.unesco.org/es/articles/preguntas-y-respuestas-por-que-es-esencial-la-educacion-digital-para-la-ciudadania-mundial}} La ausencia de dicha educación, tanto para menores como para padres y docentes, limita la capacidad de prevención.

La salud mental también desempeña un rol crítico. Investigaciones internacionales indican que adolescentes que experimentan aislamiento social, ansiedad, depresión, baja autoestima o conflictos familiares presentan mayor vulnerabilidad frente a procesos de manipulación en línea.

Desafortunadamente, tras la pandemia, el aumento en estas condiciones ha incrementado significativamente, lo cual ha empujado a muchos jóvenes a buscar validación, pertenencia o refugio en comunidades en línea.

Otro elemento clave es la falta de herramientas y conocimiento por parte de padres y tutores, dejando a los menores en una situación de desamparo digital en sus propios hogares.\footnote{UNICEF Argentina. (s. f.). Guía para el uso responsable de Internet y las redes sociales. \url{https://www.unicef.org/argentina/sites/unicef.org.argentina/files/2020-07/Guia-digital_8-julio.pdf}} La utilización de dispositivos móviles personales, combinada con aplicaciones de mensajería privada, dificulta la detección temprana de comportamientos de riesgo tanto por parte de las familias como de las autoridades.

Desde una perspectiva institucional, uno de los mayores desafíos es el carácter global de internet. Las diferencias entre legislaciones nacionales generan vacíos jurídicos que son aprovechados por organizaciones criminales. Mientras algunos países cuentan con marcos regulatorios robustos y unidades especializadas en delitos informáticos, otros carecen de recursos tecnológicos, legislación específica o personal capacitado para investigar estos delitos.

Asimismo, persiste un debate internacional respecto al equilibrio entre la protección de la infancia y la preservación de derechos fundamentales como la privacidad, la libertad de expresión y la protección de datos personales. El debate se centra en el cifrado de end-to-end encryption.\footnote{Reglamento General de Protección de Datos. (s. f.). Reglamento General de Protección de Datos (RGPD). \url{https://gdpr.eu/}} Mientras que los defensores de los derechos digitales argumentan que el cifrado protege la privacidad general, las agencias de seguridad señalan que crea “zonas ciegas” donde el abuso infantil se distribuye sin posibilidad de ser detectado por las plataformas.\footnote{Oficina del Alto Comisionado de las Naciones Unidas para los Derechos Humanos. (s. f.). Comité de los Derechos del Niño. \url{https://www.ohchr.org/en/treaty-bodies/crc}}

Las empresas tecnológicas también enfrentan importantes desafíos.\footnote{Comisión Europea. (2024). Reglamento (UE) 2024/1689 del Parlamento Europeo y del Consejo, de 13 de junio de 2024, por el que se establecen normas armonizadas en materia de inteligencia artificial (Reglamento de Inteligencia Artificial). EUR-Lex. \url{https://eur-lex.europa.eu/legal-content/EN/TXT/?uri=CELEX\%3A32024R1689}} Las plataformas deben moderar millones de publicaciones diarias mientras intentan identificar contenido ilegal mediante herramientas automatizadas y equipos humanos de revisión. La rapidez con que surgen nuevas aplicaciones, tecnologías de cifrado y métodos de evasión dificulta mantener sistemas efectivos de prevención. Los países en desarrollo también, a menudo, carecen de la infraestructura tecnológica, los recursos presupuestarios y las unidades policiales especializadas en ciberdelitos para hacer frente a estas corporaciones tecnológicas. Esta brecha tecnológica incrementa la vulnerabilidad de millones de menores y limita la capacidad de cooperación internacional.

\subsection{Posibles soluciones}

Para abordar un problema transnacional, la comunidad internacional debe evaluar enfoques multisectoriales que equilibren la soberanía estatal con la cooperación global.

Una de las estrategias más ampliamente respaldadas consiste en promover programas de alfabetización digital desde edades tempranas. Integrar la ciudadanía digital y la seguridad en línea en los currículos escolares obligatorios a nivel global, capacitando a estudiantes, profesores y familias, es el primer paso a tomar.\footnote{Organización de las Naciones Unidas para la Educación, la Ciencia y la Cultura. (s. f.). Ciudadanía digital. \url{https://www.unesco.org/es/fieldoffice/montevideo/expertise/ciudadaniadigital}} Estos programas buscan desarrollar competencias relacionadas con el pensamiento crítico, la identificación de intentos de manipulación y el uso responsable de las redes sociales. Diversos organismos internacionales sostienen que la prevención constituye la medida más eficaz para reducir la victimización infantil.

En el ámbito regulatorio, numerosos estados han comenzado a imponer mayores obligaciones a las plataformas digitales. Obligando por ley a las empresas tecnológicas a aplicar el principio de “Seguridad por Diseño”, penalizando financieramente a las plataformas que no retiren contenido criminal de forma inmediata o cuyos algoritmos promuevan la radicalización.\footnote{European Commission. (2024). Reglamento de Servicios Digitales (Digital Services Act). \url{https://digital-strategy.ec.europa.eu/en/policies/digital-services-act}} Entre las medidas discutidas se encuentran la evaluación periódica de riesgos para menores, la implementación de sistemas de verificación de edad, mecanismos ágiles para reportar contenido ilegal, mayor transparencia en los algoritmos de recomendación y obligaciones de cooperación con autoridades jurídicas.

Otra prioridad consiste en fortalecer la cooperación internacional.\footnote{WeProtect Global Alliance. (s. f.). WeProtect Global Alliance. \url{https://www.weprotect.org/}} Dado que estos delitos rara vez se limitan a un solo país, resulta indispensable mejorar el intercambio de información entre fuerzas policiales, fiscalías, organismos de protección infantil y empresas tecnológicas. Instrumentos de asistencia jurídica mutua, operaciones coordinadas e intercambio de inteligencia permiten acelerar investigaciones y rescatar víctimas.

El desarrollo de herramientas de inteligencia artificial para la detección temprana presenta otra línea de acción prometedora. Promover la expansión de dispositivos tecnológicos éticos que detecten patrones de grooming o material de abuso sin romper los principios fundamentales de privacidad de los usuarios comunes, mediante tecnologías de hashing, es esencial.\footnote{Internet Watch Foundation. (s. f.). Internet Watch Foundation. \url{https://www.iwf.org.uk/}} No obstante, estas tecnologías requieren supervisión adecuada para minimizar errores y proteger los derechos fundamentales.

De igual manera, expertos destacan la necesidad de fortalecer los mecanismos de atención integral a las víctimas. Garantizar que los países cuenten con líneas de ayuda accesibles, servicios de salud mental especializados en trauma y mecanismos ágiles para el “derecho al olvido”, asegurando que el contenido abusivo sea eliminado permanentemente de la red.\footnote{Organización Internacional de Telecomunicaciones. (s. f.). Protección de los niños en línea. \url{https://www.itu.int/en/ITU-D/Cybersecurity/Pages/COP/COP.aspx}}

\subsection{Acciones pasadas}

La comunidad internacional ha intentado construir una estructura de protección a lo largo de las últimas décadas, enfrentando tanto éxitos notables como severas limitaciones estructurales. El surgimiento de internet y la rápida evolución de las tecnologías digitales han obligado a adaptar continuamente los marcos jurídicos e institucionales existentes para responder a nuevas formas de explotación y radicalización infantil.

Si bien los primeros instrumentos internacionales fueron concebidos antes de la era digital, hoy continúan sirviendo como fundamento para elaborar políticas más específicas en relación a la protección en línea. Entre estos instrumentos jurídicos internacionales se encuentra la Convención sobre los Derechos del Niño, adoptada por la Asamblea General de las Naciones Unidas en 1989.\footnote{Organización de las Naciones Unidas. (s. f.). Día de los Derechos del Niño: Derechos del niño. \url{https://www.un.org/es/events/childrenday/pdf/derechos.pdf}} Su Protocolo facultativo relativo a la venta de niños, la prostitución infantil y la utilización de niños en pornografía infantil (2000) obliga a los Estados a criminalizar estas conductas.\footnote{Oficina del Alto Comisionado de las Naciones Unidas para los Derechos Humanos. (s. f.). Convención sobre los Derechos del Niño. \url{https://www.ohchr.org/en/instruments-mechanisms/instruments/convention-rights-child}} Aunque la convención fue redactada antes del auge del internet, sus principios han sido interpretados para aplicarse plenamente al entorno digital. Derechos como el interés superior del niño, la protección frente a la violencia, el acceso a la educación, la privacidad y la libertad de expresión deben garantizarse también en espacios virtuales.

En 2021, el Comité de los Derechos del Niño adoptó la Observación General No. 25 sobre los derechos del niño en relación con el entorno digital, uno de los documentos internacionales más importantes sobre el tema.\footnote{Oficina del Alto Comisionado de las Naciones Unidas para los Derechos Humanos. (s. f.). Comité de los Derechos del Niño. \url{https://www.ohchr.org/en/treaty-bodies/crc}} Esto reconoce expresamente que los Estados tienen la responsabilidad de garantizar que los menores puedan acceder a los beneficios del entorno digital sin exponerse a riesgos como la explotación sexual, el ciberacoso, la radicalización o manipulación algorítmica. De igual forma, invita a los gobiernos a exigir mayor responsabilidad a las empresas tecnológicas respecto de la protección infantil.

Otras iniciativas claves han sido el Fondo de las Naciones Unidas para la Infancia (UNICEF). Su papel desarrollando programas educativos, investigaciones y asistencia técnica a los gobiernos es de admirar. Su estrategia parte de un enfoque basado en los derechos del niño, buscando equilibrar la protección frente a los riesgos digitales con el acceso equitativo a la tecnología y la información.

UNICEF ha llevado a cabo campañas internacionales de alfabetización digital dirigidas tanto a menores como a padres, docentes y cuidadores, promoviendo el uso seguro de internet y la identificación temprana de situaciones de riesgo.\footnote{UNICEF. (s. f.). Cómo mantener a los niños seguros en Internet. \url{https://www.unicef.org/protection/keeping-children-safe-online}} La organización también trabaja junto con empresas tecnológicas para promover el cumplimiento de los Child Rights and Business Principles, un conjunto de principios que incentivan a las compañías a incorporar la protección infantil en el diseño de plataformas digitales, sistemas de moderación de contenido y mecanismos de denuncia.\footnote{UNICEF. (s. f.). ¿Qué hacemos? \url{https://www.unicef.org/es/que-hacemos}} Adicionalmente, UNICEF lidera campañas globales como el “We Protect Global Alliance” para concientizar y presionar a los gobiernos a implementar políticas de seguridad nacional digital.\footnote{WeProtect Global Alliance. (s. f.). Alianza Mundial WeProtect. \url{https://www.weprotect.org/}} WeProtect también ha impulsado investigaciones sobre tendencias emergentes, incluyendo el uso de la inteligencia artificial, criptomonedas y plataformas cifradas por parte de redes criminales.

Sin embargo, a pesar de su importante labor, UNICEF enfrenta limitaciones derivadas de su mandato. Como agencia de cooperación internacional, no posee facultades para imponer sanciones ni obligar a los Estados o empresas privadas a implementar determinadas medidas, dependiendo en gran medida del compromiso voluntario de sus colaboradores. Es por este motivo que la cooperación policial internacional se ha convertido en uno de los pilares fundamentales de la respuesta global.

La Organización Internacional de Policía Criminal (INTERPOL) coordina investigaciones internacionales relacionadas con delitos contra menores, facilitando el intercambio de inteligencia entre sus 196 países miembros. A través de la Base de Datos Internacional sobre la Explotación Sexual de Niños (ISCE) de INTERPOL, las fuerzas policiales comparten datos y tecnologías de reconocimiento de imágenes para identificar a las víctimas y localizar los lugares donde se producen los abusos.\footnote{INTERPOL. (s. f.). Delitos contra los niños. \url{https://www.interpol.int/Crimes/Crimes-against-children}}

A nivel regional, Europol ha fortalecido considerablemente las capacidades de los Estados miembros de la Unión Europea mediante el European Cybercrime Centre (EC3), especializado en ciberdelincuencia. Europol coordina investigaciones conjuntas, proporciona apoyo técnico a las fuerzas policiales nacionales y publica anualmente el Internet Organised Crime Threat Assessment (IOCTA), considerado uno de los principales informes sobre amenazas digitales en Europa.

No obstante, tanto INTERPOL como Europol enfrentan desafíos importantes. Su efectividad depende de la cooperación activa de los estados miembros, quienes poseen diferentes capacidades técnicas, recursos financieros y marcos legales. Además, las investigaciones suelen verse obstaculizadas por diferencias jurisdiccionales, procesos de extradición prolongados y la utilización de servicios cifrados que dificultan la obtención de evidencia digital.

Adicionalmente, la Unión Europea ha desarrollado uno de los marcos regulatorios más avanzados en materia de protección digital. El Digital Services Act (DSA) establece nuevas obligaciones para las grandes plataformas digitales.\footnote{European Commission. (2024). Reglamento de Servicios Digitales (Digital Services Act). \url{https://digital-strategy.ec.europa.eu/en/policies/digital-services-act}} Entre ellas destacan la evaluación de riesgos relacionados con menores de edad, prohíbe la publicidad dirigida basada en el perfilado de menores y establece multas multimillonarias si no protegen la salud mental y seguridad de los jóvenes. Simultáneamente, el Reglamento General de Protección de Datos (GDPR) ha fortalecido la protección de los datos personales de menores, imponiendo obligaciones adicionales a las empresas respecto al tratamiento de información infantil y al consentimiento parental en determinados servicios digitales.\footnote{Reglamento General de Protección de Datos. (s. f.). Reglamento General de Protección de Datos (RGPD). \url{https://gdpr.eu/}}
